% Generated draft: requires mathematical and source review.
\documentclass[11pt]{article}
\usepackage[margin=25mm]{geometry}
\usepackage{fontspec}
\setmainfont{Latin Modern Roman}
\usepackage{amsmath,amssymb,mathtools,unicode-math}
\setmathfont{Latin Modern Math}
\usepackage{enumitem,needspace,xurl,hyperref}
\setlength{\parindent}{0pt}
\setlength{\parskip}{5pt}
\setlength{\emergencystretch}{3em}
\newcommand{\N}{\mathbb{N}}
\newcommand{\Z}{\mathbb{Z}}
\newcommand{\Q}{\mathbb{Q}}
\newcommand{\R}{\mathbb{R}}
\newcommand{\C}{\mathbb{C}}
\newcommand{\F}{\mathbb{F}}
\newcommand{\A}{\mathbb{A}}
\newcommand{\PP}{\mathbb{P}}
\newcommand{\E}{\mathbb{E}}
\newcommand{\eps}{\varepsilon}
\newcommand{\dd}{\,\mathrm{d}}
\begin{document}
\section*{UnsolvedMath ID: 700833}
\textbf{Source problem code:} catalogue entry
\subsection{Definitions and mathematical statement}
In the conventional homogeneous interpretation of the source's phrase ``the countable universal poset,'' let $(\Omega,\leq)$ be a countably infinite partially ordered set that is universal for countable partially ordered sets and ultrahomogeneous. Universality means that every countable partially ordered set $(P,\preceq)$ admits an injective map $e:P\to\Omega$ satisfying
\[
x\preceq y\quad\Longleftrightarrow\quad e(x)\leq e(y).
\]
Ultrahomogeneity means that every isomorphism between two finite induced subposets of $\Omega$ extends to an automorphism of $(\Omega,\leq)$. The supplied record does not explicitly state ultrahomogeneity; this conventional interpretation requires confirmation against the original source.

Define
\[
G=\operatorname{Aut}(\Omega,\leq)
 =\{g:\Omega\to\Omega\mid g\text{ is bijective and }(x\leq y\Longleftrightarrow g(x)\leq g(y))\text{ for all }x,y\in\Omega\},
\]
with composition as the group operation. For a subgroup $H\leq G$, let $[G:H]$ denote the cardinality of the set of left cosets of $H$ in $G$. Say that $H$ is transitive on $\Omega$ if
\[
\forall x,y\in\Omega\ \exists h\in H\quad h(x)=y.
\]
The question is whether
\[
\forall H\leq G,\qquad
\left([G:H]<2^{\aleph_0}\ \land\ \forall x,y\in\Omega\ \exists h\in H\ (h(x)=y)\right)
\Longrightarrow H=G.
\]
Here $2^{\aleph_0}$ is the cardinality of the continuum. No topological restriction on $H$ is stated.

As background, the supplied source attributes the simplicity of $G$ to A. M. W. Glass, S. H. McCleary, and M. Rubin, in the 1993 reference listed below. This is a supplied background assertion, not a result established by this draft.
\subsection{Short English statement}
Under the conventional homogeneous interpretation of the countable universal poset, must every transitive subgroup of its automorphism group with index strictly smaller than the continuum equal the whole automorphism group? The source's intended meaning of ``countable universal poset'' requires confirmation.
\subsection{Sources}
\begin{itemize}
\item Local UnsolvedMath registry, permanent ID $700833$, exact source problem code \texttt{catalogue entry}, title ``Kourovka Notebook Problem 12.13.'' The supplied record labels the problem open; that label is not an independently verified literature assessment.
\item \emph{The Kourovka Notebook (Unsolved Problems in Group Theory)}, original problem 12.13, 12th issue (1992), proposed by A. M. W. Glass. The supplied background identifies the 21st issue edition (2026), \texttt{arXiv:1401.0300v48}; the collection description names E. I. Khukhro and V. D. Mazurov as editors.
\item A. M. W. Glass, S. H. McCleary, and M. Rubin, \emph{Math. Z.} \textbf{214}, no. 1 (1993), 55--66. The supplied statement cites this reference for simplicity of the automorphism group. An article title is not supplied.
\end{itemize}
\end{document}
